\chapter{Training und Trainingszeit-Augmentierung}
\label{kap:trainingszeit}

\section{Aufbau}
\label{sec:trainingsaufbau}

\absatz{Datenaufteilung.} Trainiert wird auf den 6000 synthetischen Bildern,
validiert auf dem \emph{realen} \code{val}-Bild (88 Objekte). Das ist eine
bewusste Entscheidung: Die Modellauswahl muss auf der Zieldomäne erfolgen, sonst
wählt man das Modell, das am besten zu den eigenen Renderartefakten passt.
\code{test} (86 Objekte) berührt Training und Modellauswahl nie und dient
ausschließlich der Endbewertung. Ein vierter Split \code{sval} aus 400
synthetischen Bildern mit anderem Seed dient allein der Messung des
Sim-to-Real-Abstands (Abschnitt~\ref{sec:simtoreal}).

\absatz{Modellwahl.} Als Baseline dient \code{yolo11s} mit 9,4~Millionen
Parametern. Die Größe folgt aus der Datenlage: Der Datensatz hat eine Klasse und
nur neun zugrunde liegende Objektinstanzen, die Formvielfalt ist klein, die
nötige Modellkapazität damit ebenfalls; die \enquote{nano}-Variante verliert bei
12--40\px{} großen Objekten Rückruf, größere Varianten würden die synthetischen
Eigenheiten eher memorieren. Ausgangspunkt sind COCO-vortrainierte Gewichte: Die
vortrainierten Low-Level-Filter sind domänenunabhängig \parencite{yosinski2014},
und \textcite{hinterstoisser2018} zeigen den Nutzen explizit für synthetische
Trainingsdaten. Die \emph{Familie} yolo11 war zunächst konservativ gewählt
(ausgereifte Implementierung, bekannte Kopfstruktur P3/P4/P5) und nicht belegt;
Abschnitt~\ref{sec:p2} prüft mit zwei YOLO26-Varianten nach, was ein feinerer
Detektionskopf bringt. Alle verglichenen Modelle haben auf P3/P4/P5 dieselben
Schrittweiten, sodass die Auflösungsargumentation und die Ergebnisse von
Abschnitt~\ref{sec:imgsz} unverändert gültig bleiben.

\absatz{Eingabegröße.} \code{imgsz}=1024, identisch zur Leinwandgröße des
Generators. Damit durchlaufen die synthetischen Bilder \emph{keine} Skalierung,
und die Objekte erreichen das Netz mit exakt der Pixelgröße, für die die Pipeline
ausgelegt ist (Anforderung~A2). Bei 640\px{} wären sie nur 7--27\px{} groß und
lägen teils unter der Auflösung der P3-Ebene.

\absatz{Optimierung.} AdamW \parencite{loshchilov2019} mit
$\text{lr}_0 = 10^{-3}$, Endlernrate 1\,\% davon, Kosinusabfall
\parencite{loshchilov2017sgdr}, drei Aufwärmepochen, Weight Decay $5\cdot10^{-4}$,
Batch~16, Mixed Precision \parencite{micikevicius2018}, fester Seed und
deterministische Kernel. AdamW ist explizit gesetzt statt \code{auto}, weil es
bei kleinen Datensätzen schneller konvergiert und robust gegenüber der
Lernratenwahl ist -- wichtig, weil kein Budget für eine Lernratensuche vorhanden
ist. Die Verlustgewichte bleiben auf den Standardwerten
(\code{box}\,7,5 / \code{cls}\,0,5): Eine Erhöhung des Boxgewichts wäre bei
kleinen Objekten denkbar, ohne separaten Abstimmungssplit aber nicht belegbar.
Das Budget beträgt 60 Epochen mit Early Stopping nach 20 Epochen ohne
Verbesserung auf dem realen \code{val}-Bild; der Hauptlauf endete dadurch nach
30 Epochen mit dem besten Modell aus Epoche~10.

\section{Augmentierungsparameter zur Trainingszeit}
\label{sec:trainparam}

Die abgeschalteten Parameter sind genauso begründungspflichtig wie die aktiven --
ein auf Null gesetzter Standardwert ist eine Entscheidung.

Aktiv sind Mosaic (immer, in den letzten zehn Epochen abgeschaltet), Skalierung
$0{,}6$--$1{,}4$, Translation, horizontale und vertikale Spiegelung sowie
HSV-Variation (Ton 0,015, Sättigung 0,7, Helligkeit 0,4). Abgeschaltet sind
\code{degrees}, \code{shear} und \code{perspective}, weil Bildrotation
achsparallele Boxen aufbläht -- bei $10^\circ$ um 15,8\,\% der Kantenlänge, was
die Überlappung mit der korrekten Box auf 0,745 senkt
(Abschnitt~\ref{sec:boxrotation}) --, sowie \code{mixup} und \code{cutmix}, die
Bildinhalte bei erhaltenen Labels überlagern und damit Geisterobjekte mit vollem
Boxlabel erzeugen \parencite{zhang2018mixup,yun2019cutmix}: Bei dicht gepackten
20\px-Objekten überwiegt der Labelkonflikt. \code{erasing} entfällt, weil es in
der Werkzeugkette nur bei der Klassifikation wirkt, und \code{copy\_paste},
weil es Segmentierungsmasken benötigt -- beides übernimmt der Generator
physikalisch korrekt. Die vollständige Tabelle mit Begründung je Parameter steht
im Anhang (Tabelle~\ref{tab:trainaug}).

Die beiden Augmentierungsschichten wirken bei der Objektgröße multiplikativ: Der
Generator erzeugt $[12, 44]$\px, die Trainingsskalierung multipliziert mit
$[0{,}6;\,1{,}4]$, sodass das Netz effektiv Objekte von
$12 \cdot 0{,}6 = 7{,}2$ bis $44 \cdot 1{,}4 = 61{,}6$\px{} sieht. Diese Spanne
deckt nach unten die Objektgrößen bei niedriger
Inferenzauflösung und nach oben jene bei starker Hochskalierung ab -- also den
gesamten in Abschnitt~\ref{sec:imgsz} durchmusterten Bereich. Der untere Rand von
7,2\px{} liegt knapp unter dem Stride der feinsten Ebene; ein kleinerer Wert für
\code{scale} hätte den Bereich sinnvoll auflösbarer Objekte verlassen.

\section{Trainingsverlauf}
\label{sec:trainingsverlauf}


Die Lernkurven zeigen den Sim-to-Real-Abstand in zeitlicher Form: Der
Trainingsverlust fällt über alle 30 Epochen monoton, die Metrik auf dem realen
\code{val}-Bild ist ab Epoche~10 flach. Ab dort verbessert das
Modell ausschließlich seine Anpassung an Eigenschaften der synthetischen Daten,
die in der Zieldomäne nicht gelten. Genau deshalb ist es methodisch notwendig,
das beste Modell über den \emph{realen} Validierungssplit auszuwählen: Eine
synthetische Validierung wählte das Modell aus der letzten Epoche und damit das
am stärksten überangepasste.
